% =====================================================================
\chapter{Descenso de gradiente}
\label{tema:2}
% =====================================================================

El descenso de gradiente es el algoritmo del que derivan todos los optimizadores del módulo. Este tema demuestra sus garantías de convergencia en tres situaciones de generalidad decreciente: funciones suaves arbitrarias, funciones fuertemente convexas y funciones cuadráticas, en las que el análisis es exacto. El resultado central es que el número de iteraciones crece proporcionalmente al número de condición $\kappa$. Se sigue~\cite[§9.2 y §9.3]{bv}, \cite[§12.3]{d2l} y las clases~21 y~22 de~\cite{mit18065}.

\section{El método y el lema de descenso}

\begin{definicion}[Descenso de gradiente]\label{def:gd}
Dados $\theta_0\in\R^{n}$ y una \emph{tasa de aprendizaje} (o \emph{paso}) $\eta>0$, el \emph{descenso de gradiente} genera la sucesión
\[
\theta_{k+1}=\theta_k-\eta\,\nabla f(\theta_k),\qquad k=0,1,2,\dots
\]
\end{definicion}

La dirección $-\nabla f(\theta_k)$ es la de máximo descenso local (módulo~4, tema~1). Por el desarrollo de Taylor de primer orden, $f(\theta_k-\eta\nabla f(\theta_k))=f(\theta_k)-\eta\norm{\nabla f(\theta_k)}^{2}+o(\eta)$, de modo que un paso suficientemente pequeño reduce $f$ siempre que el gradiente no sea nulo. El lema siguiente cuantifica «suficientemente pequeño».

\begin{lema}[Lema de descenso]\label{lem:descenso}
Si $f$ es $L$-suave, para todo $x,y\in\R^{n}$,
\[
f(y)\le f(x)+\nabla f(x)\T(y-x)+\frac{L}{2}\norm{y-x}^{2}.
\]
\end{lema}

\begin{corolario}[Descenso suficiente]\label{cor:descenso}
Si $f$ es $L$-suave y $\theta^{+}=\theta-\eta\nabla f(\theta)$ con $0<\eta<2/L$, entonces
\[
f(\theta^{+})\le f(\theta)-\eta\Bigl(1-\frac{L\eta}{2}\Bigr)\norm{\nabla f(\theta)}^{2}.
\]
Con $\eta=1/L$ la cota es $f(\theta^{+})\le f(\theta)-\frac{1}{2L}\norm{\nabla f(\theta)}^{2}$, y es la mejor que da el lema.
\end{corolario}

\begin{teorema}[Caso no convexo]\label{teo:gd-noconvexo}
Sea $f$ $L$-suave y acotada inferiormente por $f^{\ast}$. El descenso de gradiente con $\eta=1/L$ cumple, para todo $K\ge 1$,
\[
\min_{0\le k<K}\norm{\nabla f(\theta_k)}^{2}\le\frac{2L\bigl(f(\theta_0)-f^{\ast}\bigr)}{K}.
\]
\end{teorema}

El \cref{teo:gd-noconvexo} solo garantiza aproximarse a un punto crítico, que puede ser una silla; es la clase de garantía que cabe esperar para la pérdida de una red profunda.

\section{Convergencia lineal en funciones fuertemente convexas}

\begin{teorema}[Convergencia lineal]\label{teo:gd-fuerte}
Sea $f$ $\mu$-fuertemente convexa y $L$-suave, con valor óptimo $f^{\ast}$. El descenso de gradiente con $\eta=1/L$ cumple
\[
f(\theta_k)-f^{\ast}\le\Bigl(1-\frac{1}{\kappa}\Bigr)^{k}\bigl(f(\theta_0)-f^{\ast}\bigr),\qquad\kappa=\frac{L}{\mu}.
\]
\end{teorema}

\begin{corolario}[Número de iteraciones]\label{cor:iteraciones}
En las condiciones del \cref{teo:gd-fuerte}, para alcanzar $f(\theta_k)-f^{\ast}\le\epsilon\,(f(\theta_0)-f^{\ast})$ basta con $k\ge\kappa\ln(1/\epsilon)$ iteraciones.
\end{corolario}

La convergencia se llama \emph{lineal} porque el error se multiplica en cada paso por un factor constante menor que~1: el logaritmo del error decrece linealmente con $k$. El factor depende solo de $\kappa$. El análisis de~\cite[§9.3]{bv} obtiene el mismo factor $1-\mu/L$, salvo constantes, para búsquedas lineales exacta y con retroceso.

\section{Análisis exacto en cuadráticas}

En una función cuadrática el método es lineal en $\theta$ y puede resolverse en forma cerrada. Toda función de clase $C^{2}$ es aproximadamente cuadrática cerca de un mínimo no degenerado, de modo que este caso describe el comportamiento asintótico general.

\addtocounter{definicion}{1}%

\begin{teorema}[Descenso de gradiente en cuadráticas]\label{teo:gd-cuadratica}
Sea $f(\theta)=\tfrac12\theta\T Q\theta$ con $Q$ simétrica de autovalores $0<\mu=\lambda_1\le\dots\le\lambda_n=L$. Sea $\rho(\eta)=\max_i\abs{1-\eta\lambda_i}$.
\begin{enumerate}
  \item $\norm{\theta_k}\le\rho(\eta)^{k}\norm{\theta_0}$, con igualdad si $\theta_0$ es autovector del autovalor que realiza el máximo.
  \item El método converge a $\theta^{\ast}=0$ para todo $\theta_0$ si y solo si $0<\eta<2/L$.
  \item El valor de $\eta$ que minimiza $\rho(\eta)$ es $\eta^{\ast}=2/(\mu+L)$, y
  \[
  \rho(\eta^{\ast})=\frac{L-\mu}{L+\mu}=\frac{\kappa-1}{\kappa+1}.
  \]
\end{enumerate}
\end{teorema}

\begin{observacion}
El \cref{teo:gd-cuadratica} explica la conducta del método. La dirección de curvatura máxima, $\lambda_n=L$, limita el paso: $\eta<2/L$. La dirección de curvatura mínima, $\lambda_1=\mu$, se contrae entonces con un factor $1-\eta\mu>1-2/\kappa$, próximo a~1 si $\kappa$ es grande. Con el paso óptimo, las dos direcciones extremas se contraen al mismo ritmo, pero con signos opuestos: la de curvatura alta oscila.
\end{observacion}

\addtocounter{definicion}{1}%

\addtocounter{figure}{1}%

\begin{corolario}\label{cor:iter-cuadratica}
En las condiciones del \cref{teo:gd-cuadratica}, con el paso óptimo, el menor $k$ con $\rho^{k}\le\epsilon$ es $k=\lceil\ln\epsilon/\ln\rho\rceil$, que para $\kappa$ grande es aproximadamente $\frac{\kappa}{2}\ln(1/\epsilon)$.
\end{corolario}

El \cref{tab:kappa} aplica el corolario con $\epsilon=10^{-3}$. La última columna anticipa el resultado del \cref{tema:4}: con momentum, el número de iteraciones crece como $\sqrt\kappa$ en lugar de como $\kappa$.

\begin{table}[h]
\centering\small
\caption{Factor de contracción e iteraciones necesarias para reducir el error a $10^{-3}$ en una cuadrática, con parámetros óptimos.}\label{tab:kappa}
\begin{tabular}{@{}rcrcr@{}}
\toprule
$\kappa$ & $\rho$ del gradiente & Iteraciones & $\rho$ con momentum (\cref{teo:heavy-ball}) & Iteraciones \\
\midrule
10 & $0{,}8182$ & 35 & $0{,}5195$ & 11 \\
100 & $0{,}9802$ & 346 & $0{,}8182$ & 35 \\
1000 & $0{,}9980$ & 3454 & $0{,}9387$ & 110 \\
\bottomrule
\end{tabular}
\end{table}

\section{La función de Rosenbrock y el borde de estabilidad}

\addtocounter{definicion}{1}%

\begin{observacion}[Borde de estabilidad]\label{obs:borde}
La condición $\eta<2/\lambda_{\max}$ también se observa en redes profundas, pero en sentido inverso. Cohen \emph{et al.}~\cite{cohen} observaron que, al entrenar con descenso de gradiente de paso fijo, la curvatura máxima $\lambda_{\max}$ del hessiano de la pérdida crece hasta situarse en torno a $2/\eta$ y permanece ahí; la pérdida deja de decrecer de forma monótona a corto plazo, pero sigue bajando a largo plazo. La red evoluciona hacia el límite de estabilidad del \cref{teo:gd-cuadratica}.
\end{observacion}

